\chapter{Birational geometry and alterations}
\index{birational geometry}\index{alteration}\index{flattening}
\label{app:birational-geometry}

This appendix records the birational constructions used by the model and
envelope theories.  The guiding distinction is between a
\emph{modification}, which leaves the function field unchanged, and an
\emph{alteration}, which is only generically finite.  Resolution gives the
first in characteristic zero; de Jong's theorem gives the second in
arbitrary characteristic over a perfect field.

All varieties in this appendix are integral separated schemes of finite
type over a field unless a different hypothesis is stated.  A birational
morphism is always understood to induce the chosen identification of
function fields.

\section{Birational morphisms and modifications}
\label{sec:appB-modifications}

\begin{definition}[Modification]
\label{def:appB-modification}
A modification of an integral scheme \(X\) is a proper birational morphism
\(\mu:Y\to X\).  It is projective if \(\mu\) is projective.  The exceptional
locus is the complement of the largest open set over which \(\mu\) is an
isomorphism.
\end{definition}

\begin{lemma}\label{lem:appB-common-domination}
Let \(Y_1\to X\) and \(Y_2\to X\) be projective modifications.  There is a
projective modification \(Z\to X\) dominating both.  If \(X,Y_1,Y_2\) are
normal, \(Z\) may be chosen normal.
\end{lemma}

\begin{proof}
The rational map \(Y_1\dashrightarrow Y_2\) is defined over the common
dense open on which both maps are isomorphisms.  Take the schematic closure
of its graph in \(Y_1\times_XY_2\).  The two projections are projective and
birational.  The component containing the graph is integral and dominates
\(X\).  Normalize it for the final assertion; normalization is finite
because the schemes are of finite type over a field.
\end{proof}

\begin{remark}[Checkpoint]
\label{rem:appB-model-nets}
The family of models is directed by domination, not usually by a countable
sequence.  Limits over all models in
\Cref{chap:models,chap:bj-homogeneity-bdivisors} are therefore nets unless a separate
countability hypothesis has been imposed.
\end{remark}

\section{Blowups and exceptional divisors}
\label{sec:appB-blowups}

Let \(\mathcal I\subset\cO_X\) be a nonzero coherent ideal and
\(\pi:Y=\operatorname{Bl}_{\mathcal I}X\to X\) its blowup.  By
\Cref{prop:appA-blowup-proj},
\(\mathcal I\cO_Y\) is invertible.

\begin{definition}[Exceptional Cartier divisor]
\label{def:appB-exceptional-divisor}
The effective Cartier divisor \(E\) determined by
\[
 \mathcal I\cO_Y=\cO_Y(-E)
\tag{B.1}\label{eq:appB-exceptional-ideal}
\]
is the exceptional divisor of the ideal.  The relatively ample
tautological bundle is \(\cO_Y(-E)=\cO_Y(1)\).
\end{definition}

\begin{remark}[Sign checkpoint]
\label{rem:appB-blowup-sign}
The effective divisor is \(E\), while the ideal itself is
\(\cO_Y(-E)\).  Under the metric convention of
\Cref{app:normalizations},
\(\log|s_E|(v)=-v(E)\).  Confusing \(E\) with \(-E\) reverses the PL
function attached to a flag ideal.
\end{remark}

\begin{example}[Blowing up the plane]
\label{ex:appB-blowup-plane}
For \(X=\Spec k[x,y]\) and \(\mathcal I=(x,y)\), the blowup is covered by
\[
 U_x=\Spec k[x,y/x],
 \qquad
 U_y=\Spec k[x/y,y].
\]
On \(U_x\), the exceptional divisor is \(x=0\); on \(U_y\), it is
\(y=0\).  Their transition relation \(y=x(y/x)\) shows that they glue to
\(E\simeq\mathbf P^1_k\).  The pullback of \((x,y)\) is generated by
\(x\) on \(U_x\) and by \(y\) on \(U_y\), verifying
\eqref{eq:appB-exceptional-ideal}.
\end{example}

\begin{proposition}\label{prop:appB-blowup-base-change}
For any flat morphism \(X'\to X\),
\[
 \operatorname{Bl}_{\mathcal I\cO_{X'}}X'
 \simeq
 \operatorname{Bl}_{\mathcal I}X\times_XX'.
\]
For arbitrary base change this identity can fail.
\end{proposition}

\begin{proof}
Flatness gives
\[
 \mathcal I^m\otimes_{\cO_X}\cO_{X'}
   \simeq(\mathcal I\cO_{X'})^m.
\]
Thus the Rees algebra commutes with base change, and relative Proj does as
well.  Without flatness, tensoring need not preserve the image defining
\(\mathcal I^m\).
\end{proof}

\begin{definition}[Strict transform]
\label{def:appB-strict-transform}
Let \(Z\subset X\) be a closed integral subscheme not contained in
\(V(\mathcal I)\).  Its strict transform in
\(\operatorname{Bl}_{\mathcal I}X\) is the closure of
\[
 Z\setminus V(\mathcal I)\longrightarrow
 X\setminus V(\mathcal I)
\]
inside the blowup.
\end{definition}

\section{Normalized blowups and Rees divisors}
\label{sec:appB-normalized-blowups}

Assume \(X\) is normal and \(\mathcal I\ne0\).  Let
\[
 \pi^\nu:Y^\nu\longrightarrow X
\tag{B.2}\label{eq:appB-normalized-blowup}
\]
be the normalized blowup.  The prime components
\(E_1,\ldots,E_r\) of the effective Cartier divisor defined by
\(\mathcal I\cO_{Y^\nu}\) determine divisorial valuations
\(\ord_{E_i}\) on \(K(X)\).

\begin{theorem}\label{thm:appB-rees-valuations}
There are positive integers
\(a_i=\ord_{E_i}(\mathcal I)\) such that, for all \(m\geq1\),
\[
 \overline{\mathcal I^m}
 =\{f\in\cO_X:
       \ord_{E_i}(f)\geq ma_i\text{ for }1\leq i\leq r\}.
\tag{B.3}\label{eq:appB-rees-valuations}
\]
The assertion is local on \(X\), and an inequality for a section means the
corresponding inequality in every local ring where it is defined.
\end{theorem}

\begin{proof}
On \(Y^\nu\), the extended ideal is
\(\cO_{Y^\nu}(-E)\), with
\(E=\sum_i a_iE_i\).  Some \(E_i\) may dominate a divisor on \(X\);
Rees divisors are not necessarily exceptional.  A rational function \(f\)
belongs to the pushforward
\(\pi^\nu_*\cO_{Y^\nu}(-mE)\) precisely when its divisor has coefficient at
least \(ma_i\) along each \(E_i\), together with regularity away from the
exceptional locus.  Normality tests regularity in codimension one.
For completeness, its identification with
\(\overline{\mathcal I^m}\) can be checked without a saturation claim
about graded pieces.  If \(f\) is integral over \(\mathcal I^m\), divide
its integral equation on a chart by a generator \(g^m\) of
\(\mathcal I^m\cO_{Y^\nu}\).  The quotient \(f/g^m\) is integral over
the normal chart ring, hence regular, so \(f\) is in the pushforward.
Conversely, every valuation ring of \(K(X)\) containing a given affine
base ring lifts to the blowup, since its extended finite-type ideal is
principal.  It also lifts to the normalization, since a valuation ring
is integrally closed.  Pulling a section of \(\cO_{Y^\nu}(-mE)\) back
to this ring gives \(v(f)\geq m v(\mathcal I)\).  The valuative
criterion in \Cref{thm:appA-valuative-integral-closure} proves the reverse
inclusion.  This proves \eqref{eq:appB-rees-valuations}.
\end{proof}

\begin{example}[Rees divisor of a monomial ideal]
\label{ex:appB-monomial-rees}
For \(\mathfrak a=(x^2,y^3)\), the compact edge of the Newton polygon has
primitive inward normal \((3,2)\).  The corresponding Rees valuation has
\[
 v(x)=3,\qquad v(y)=2,\qquad v(\mathfrak a)=6.
\]
Its inequality \(3i+2j\geq6m\) cuts out the exponent set of
\(\overline{\mathfrak a^m}\).
\end{example}

\section{Dominating a modification by a blowup}
\label{sec:appB-domination-blowup}

\begin{theorem}\label{thm:appB-admissible-domination}
Let \(S\) be quasi-compact and quasi-separated, and let
\(f:Y\to S\) be proper.  Suppose \(f\) is an isomorphism over a
quasi-compact open \(U\subset S\).  Then some blowup
\(S'\to S\), whose center is disjoint from \(U\), factors through \(Y\).
\end{theorem}

\begin{proof}[Foundational input]
We import exactly this domination statement from
\cite[\href{https://stacks.math.columbia.edu/tag/081T}{Tag 081T}]
{stacks-project}.  It produces domination by a blowup; it does not assert
flatness of any sheaf.
\end{proof}

This theorem explains why Cartier b-divisors may be studied on blowup
models.  It is first used in the model-cofinality argument of
\Cref{chap:models}, and again when realizing rational PL data in
\Cref{chap:bj-cartier-realization}.

\section{Flattening by blowup}
\label{sec:appB-flattening}

\begin{definition}[Pure transform]
\label{def:appB-pure-transform}
Let \(f:X\to S\), let \(U\subset S\) be open, and let
\(\mathcal F\) be a quasi-coherent \(\cO_X\)-module.  For a modification
\(p:S'\to S\) that is an isomorphism over \(U\), write
\[
 X'=X\times_SS',\qquad
 j:X'_{U}\hookrightarrow X'.
\]
The pure transform \(\mathcal F'\) is the image
\[
 p_X^*\mathcal F\longrightarrow
 j_*(p_X^*\mathcal F|_{X'_U}).
\tag{B.4}\label{eq:appB-pure-transform}
\]
It is the pullback modulo sections supported away from \(U\).
\end{definition}

\begin{theorem}\label{thm:appB-flattening}
Let \(S\) be quasi-compact and quasi-separated, let
\(f:X\to S\) be quasi-compact and locally of finite presentation, and let
\(\mathcal F\) be a finite-type quasi-coherent module.  Suppose that over a
quasi-compact open \(U\subset S\), the restriction \(\mathcal F_U\) is
finitely presented and flat over \(U\).  Then there is a \(U\)-admissible
blowup \(S'\to S\) for which the pure transform \(\mathcal F'\) is finitely
presented and flat over \(S'\).
\end{theorem}

\begin{proof}[Foundational input and exact provenance]
The theorem is the flattening technique of Raynaud and Gruson
\cite{RaynaudGruson1971}, in the form
\cite[\href{https://stacks.math.columbia.edu/tag/0815}{Tag 0815}]
{stacks-project}.  Its proof constructs a universal flattening center by
successively controlling finite presentations.  We import that
construction as a deep theorem.
\end{proof}

\begin{corollary}\label{cor:appB-flatify-finite-algebra}
Let \(h:T\to X\) be finite dominant between integral noetherian schemes.
There is a projective birational \(\beta:Z\to X\) such that the pure
transform of \(h_*\cO_T\) is a finite locally free \(\cO_Z\)-algebra of
rank \([K(T):K(X)]\).
\end{corollary}

\begin{proof}
Generic freeness gives a dense open \(U\) over which \(h_*\cO_T\) is
locally free of the generic rank.  Apply
\Cref{thm:appB-flattening}.  The kernel removed in
\eqref{eq:appB-pure-transform} is a two-sided ideal, so the pure transform
retains an algebra structure.  A finite, flat, finitely presented module is
locally free; see
\cite[\href{https://stacks.math.columbia.edu/tag/02K9}{Tag 02K9}]
{stacks-project}.  Its positive generic rank is constant on the integral
base.
\end{proof}

\begin{remark}[Checkpoint]
\label{rem:appB-flatification-new-source}
Flattening does not make the original map \(T\to X\) flat.  It constructs
a new finite flat map after changing both source and target.  In
\Cref{chap:bj-finite-descent}, the source is
\(\Spec_Z\mathcal A'\), and its map back to \(T\) is tracked explicitly.
\end{remark}

\section{Stein factorization}
\label{sec:appB-stein}

\begin{theorem}\label{thm:appB-stein-factorization}
Let \(f:Y\to X\) be proper, with \(X\) locally noetherian.  Then
\[
 Y\xrightarrow{g}T\xrightarrow{h}X,
\qquad
 T=\Spec_X(f_*\cO_Y),
\tag{B.5}\label{eq:appB-stein}
\]
where \(g\) is proper with geometrically connected fibres,
\(g_*\cO_Y=\cO_T\), and \(h\) is finite.
\end{theorem}

\begin{proof}[Foundational input]
This is exactly the Noetherian Stein factorization statement in
\cite[\href{https://stacks.math.columbia.edu/tag/03H0}{Tag 03H0}]
{stacks-project}.  The relative spectrum is finite because
\(f_*\cO_Y\) is coherent under the stated proper noetherian hypotheses.
Connectedness of geometric fibres is the substantive Stein assertion.
\end{proof}

\begin{proposition}\label{prop:appB-stein-normal}
In \eqref{eq:appB-stein}, assume that \(Y\) is normal and integral and that
every quasi-compact open of \(Y\) has finitely many irreducible components.
Then \(T\) is normal and integral.
\end{proposition}

\begin{proof}
The Stein middle term is the normalization of \(X\) in \(Y\).
The normality and component statement is
\cite[\href{https://stacks.math.columbia.edu/tag/035L}{Tag 035L}]
{stacks-project}.  Since \(Y\) has one component, so does \(T\).
\end{proof}

\begin{lemma}\label{lem:appB-stein-birational}
If \(f:Y\to X\) in \eqref{eq:appB-stein} is dominant and generically
finite, then \(K(T)=K(Y)\), so \(g:Y\to T\) is birational.
\end{lemma}

\begin{proof}
Over the generic point \(\eta\) of \(X\), the fibre
\(Y_\eta=\Spec K(Y)\) is finite.  Its ring of global functions is
\(K(Y)\).  Therefore the generic fibre of
\(T=\Spec_Xf_*\cO_Y\) is \(\Spec K(Y)\), giving \(K(T)=K(Y)\).
\end{proof}

\begin{example}[Why Stein is useful for alterations]
\label{ex:appB-stein-alteration}
A projective alteration \(Y\to X\) need not be birational.  Its Stein factorization
separates it into
\[
 \underbrace{Y\to T}_{\text{projective birational}}
 \quad\text{and}\quad
 \underbrace{T\to X}_{\text{finite}}.
\]
The envelope proof uses birational reflection on the first arrow and
determinant-norm descent on the second.
\end{example}

\section{Resolution in characteristic zero}
\label{sec:appB-resolution}

\begin{definition}[Log resolution]
\label{def:appB-log-resolution}
Let \(X\) be smooth over a characteristic-zero field and
\(\mathfrak a_1,\ldots,\mathfrak a_r\) nonzero coherent ideals.  A log
resolution is a projective birational morphism \(\mu:Y\to X\), with \(Y\)
smooth, such that
\[
 \mathfrak a_i\cO_Y=\cO_Y(-F_i)
\]
and the union of the exceptional locus with
\(\sum_i\operatorname{Supp}F_i\) is a simple normal crossings divisor.
\end{definition}

\begin{theorem}\label{thm:appB-resolution}
Over a field of characteristic zero, every integral variety admits a
projective resolution of singularities.  More generally, finitely many
nonzero coherent ideals on a smooth variety admit a simultaneous log resolution,
which may be chosen to be an isomorphism where the variety is smooth and
the ideals already define a simple normal crossings divisor.
\end{theorem}

\begin{proof}[Foundational input and scope]
Resolution is Hironaka's theorem
\cite{Hironaka1964I,Hironaka1964II}; the stated principalization form
follows by applying
embedded resolution to the product of the ideals.  We use only varieties
of finite type over a characteristic-zero field.  No resolution theorem is
invoked in positive or mixed characteristic.
\end{proof}

\begin{remark}[Analytic form]
\label{rem:appB-analytic-resolution}
For reduced complex analytic spaces, local resolution and principalization
are available in the analytic category; see \cite{Wlo09}.  This analytic
statement is used in \cref{part:complex}, whereas multiplier ideals of algebraic ideals
in \Cref{app:ideals} use algebraic log resolutions.
\end{remark}

\section{Alterations in arbitrary characteristic}
\label{sec:appB-alterations}

\begin{definition}[Alteration]
\label{def:appB-alteration}
An alteration is a proper, dominant, generically finite morphism
\[
 f:Y\longrightarrow X
\tag{B.6}\label{eq:appB-alteration}
\]
between integral varieties.  Its degree is \([K(Y):K(X)]\).
\end{definition}

\begin{theorem}\label{thm:appB-de-jong}
Let \(X\) be an integral variety over a perfect field \(k\).  There exists
a projective alteration \(f:Y\to X\) with \(Y\) smooth over \(k\).
If \(X\) is projective, \(Y\) may be taken projective.
\end{theorem}

\begin{proof}[Foundational input and exact provenance]
This is \cite[Theorem~4.1]{deJong1996}.  Projectivity of the alteration
morphism follows as follows.  The source is quasi-projective over \(k\);
after choosing a projective compactification \(\overline Y\), the graph
\((j,f):Y\to\overline Y\times_kX\) is a locally closed immersion.  It is
proper over \(X\), since its source is proper and its target is separated
over \(X\), and hence it is a closed immersion.
Consequently \(f\) is projective.  If \(X\) is projective, then \(Y\) is
proper over \(k\); the open immersion of \(Y\) into the regular projective
compactification furnished in the cited proof is therefore also proper,
so its image is open and closed.  On the component dominating \(X\) it is
an equality, and \(Y\) is projective over \(k\).  Perfectness ensures that
regular finite-type \(k\)-schemes are smooth.  Over an imperfect field one
must distinguish regularity from smoothness.
\end{proof}

\begin{example}[Alteration degree matters]
\label{ex:appB-frobenius-alteration}
Let \(X\) be a smooth \(n\)-dimensional variety over a perfect field \(k\) of
characteristic \(p\).  The relative Frobenius
\(F_{X/k}:X\to X^{(p)}\) is a finite surjective radicial \(k\)-morphism of
degree \(p^n\).  It is an alteration, not a resolution: it is not birational
in positive dimension.  The absolute Frobenius has the same underlying
universal-homeomorphism behavior but need not be a \(k\)-morphism unless the
Frobenius of \(k\) is the identity.  Any pushforward or norm formula must
retain the degree \(p^n\).
\end{example}

\begin{proposition}\label{prop:appB-alteration-stein-ledger}
Let \(f:Y\to X\) be the smooth alteration with projective morphism supplied by
\Cref{thm:appB-de-jong}, with \(X\) normal.  Then
\[
 Y\xrightarrow{g}T\xrightarrow{h}X
\]
has the following properties:
\begin{enumerate}
\item \(T\) is normal and integral; if \(X\) is projective over \(k\),
      then so is \(T\);
\item \(g\) is projective birational with connected fibres;
\item \(h\) is finite of degree \([K(Y):K(X)]\).
\end{enumerate}
\par
\end{proposition}

\begin{proof}
Apply \Cref{thm:appB-stein-factorization}.  Normality and integrality
come from \Cref{prop:appB-stein-normal}; birationality comes from
\Cref{lem:appB-stein-birational}.  Finiteness is part of Stein
factorization, so \(h\) is projective.  Since \(Y\to X\) is projective
and \(T\to X\) is separated, the same graph argument used in
\Cref{thm:appB-de-jong} makes \(g\) projective.  Finally, when \(X\) is
projective over \(k\), finiteness of \(h\) makes \(T\) projective over
\(k\).
\end{proof}

\section{What the envelope proof imports}
\label{sec:appB-envelope-imports}

The proof in \Crefrange{chap:bj-finite-descent}{chap:bj-alterations}
uses the following chain, with no hidden resolution assumption:
\[
\begin{tikzcd}[column sep=large]
Y \arrow[r,"g"] &
T \arrow[r,"h"] &
X\\
\text{smooth} \arrow[r,phantom,"\text{birational}" description] &
\text{normal} \arrow[r,phantom,"\text{finite}" description] &
\text{normal}.
\end{tikzcd}
\tag{B.7}\label{eq:appB-envelope-chain}
\]
The smooth theorem is applied on \(Y\).  Birational reflection transfers
the envelope property from \(Y\) to \(T\).  Finite descent transfers it
from \(T\) to \(X\).  The latter first flatifies \(h_*\cO_T\), then applies
the determinant norm.  Neither step asserts that psh functions descend
pointwise through an arbitrary alteration.

\begin{remark}[Hypothesis audit]
\label{rem:appB-envelope-hypotheses}
\begin{enumerate}
\item Integrality supplies one function field and a well-defined degree.
\item Projectivity supplies properness and keeps all models projective.
\item Perfectness enters only when the regular alteration source is called
      smooth.
\item Normality of \(X\) is used at the normal-case stage.  The later
      passage from normalization to a unibranch target uses universal
      homeomorphisms and \Cref{app:valuation-algebra}, not
      alterations alone.
\end{enumerate}
\end{remark}

\section{Dependency map}
\label{sec:appB-dependency-map}

\begin{center}
\begin{tabular}{p{0.35\textwidth}p{0.50\textwidth}}
\toprule
Birational tool & First substantial use\\
\midrule
Blowup and exceptional divisor
 & model metrics, \Cref{chap:models}\\
Normalized blowup and Rees valuations
 & test configurations, \Cref{chap:bj-test-configurations}\\
Common domination
 & b-divisors, \Cref{chap:bj-homogeneity-bdivisors}\\
Flattening by blowup
 & finite descent, \Cref{chap:bj-finite-descent}\\
Stein factorization
 & normal reduction, \Cref{chap:bj-alterations}\\
Log resolution
 & multiplier ideals, \Cref{app:ideals}\\
Smooth alteration
 & arbitrary-characteristic normal case,
   \Cref{chap:bj-alterations}\\
\bottomrule
\end{tabular}
\end{center}

\section{Source note}
\label{sec:appB-source-note}
Blowup and normalization facts are cross-audited in
\Cref{app:stacks}.  The flattening theorem is due to Raynaud and
Gruson \cite{RaynaudGruson1971}; the form used here is Tag 0815.
Resolution is \cite{Hironaka1964I,Hironaka1964II}, and alterations are
\cite{deJong1996}.  The exact Stein factorization and relative-normality
hypotheses are Tags 03H0 and 035L.
